\begin{table}[!htbp]
\centering
\footnotesize
\caption{Snapshot-specific canonical MemoryDecay parameters and model comparison.}
\label{tab:si-memorydecay-fits}
\renewcommand{\arraystretch}{1.05}
\begin{tabular}{lllllllll}
\toprule
Snapshot & $N$ & $p$ & $r$ & $q$ & \makecell[b]{Near-coincident\\ rates} & \makecell[b]{RMSE$_{\log}$\\ canonical} & \makecell[b]{RMSE$_{\log}$\\ exponential} & \makecell[b]{RMSE$_{\log}$\\ log-normal} \\
\midrule
October 2016 & 85.20 & 0.2743 & 0.5354 & 0.0249 & no & 0.0515 & 0.0664 & 0.1259 \\
July 2017 & 84.88 & 0.2123 & 0.5232 & 0.0235 & no & 0.0659 & 0.0731 & 0.1517 \\
August 2022 & 66.38 & 0.0136 & 0.0185 & 0.0321 & yes & 0.0593 & 0.0770 & 0.1441 \\
August 2025 & 69.40 & 0.0119 & 0.0221 & 0.0340 & yes & 0.0868 & 0.1133 & 0.1954 \\
\bottomrule
\end{tabular}
\begin{flushleft}\footnotesize
Canonical fits use the log-response objective, the MemoryDecay starting grid, and early-age weighting on never-film age means. The near-coincident-rates flag identifies fits with $p+r-q<0.005$; decay-equivalent diagnostics include these curves. Exponential and log-normal comparison models are fitted to the same age means with the same weights; the log-normal comparison is the only model using log(age).
\end{flushleft}
\end{table}
